\Needspace{12\baselineskip}
\section{References and Source Boundary}

\noindent\emph{These sources materially govern the article. Exact loci distinguish doctrine, theological articulation, spiritual witness, historical evidence, and editorial synthesis. Linked public translations are working study texts rather than claims to critical-edition status. Online evidence was checked through 13 July 2026.}

% Keep the audit trail readable without allowing short bibliography fragments
% to spill onto a page of their own.
\begingroup
\setlength{\parskip}{0.15em}
\titlespacing*{\subsection}{0pt}{1.35ex}{0.45ex}

\subsection*{Sacred Scripture}

{\footnotesize
\begin{itemize}[itemsep=0.1em,topsep=0.15em]
\item Gen 50:15--21; Deut 30:15--20; Sir 15:11--20; Wis 8:1, 11:20--26, 12:13--18; Job 1:20--22, 2:9--10, 42:2; Pss 22(23), 113B(115):3, 134(135):6; Prov 16:9, 19:21; Matt 6:9--34, 10:28--31, 22:34--40, 26:36--46; Luke 10:25--37, 13:1--5, 22:39--46; John 9:1--3, 13:34--35, 15:1--17, 19:10--11; Acts 2:22--24, 4:27--28, 17:24--31.
\item Rom 5--8, especially 8:18--39; 1 Cor 10:13 and 13; Gal 5:1, 13--26; Eph 1:11, 2:8--10; Phil 2:12--13; 1 Thess 5:16--18; Heb 10:5--10; Jas 1:12--18, 4:13--17; 1 Pet 4:19, 5:6--7; 1 John 3:11--24, 4:7--21.
\item Direct phrases are concise working renderings checked at the named loci; no extended modern Bible translation is reproduced.
\end{itemize}}

\subsection*{Councils, Magisterium, and authoritative synthesis}

{\footnotesize
\begin{itemize}[itemsep=0.1em,topsep=0.15em]
\item First Vatican Council, dogmatic constitution \emph{Dei Filius} (1870), ch.~1, DS 3001--3003: the one simple God is omnipotent and infinite in intellect and will; Providence encompasses all things, including events arising from creatures' free action; \sourceurl{https://www.vatican.va/archive/hist_councils/i-vatican-council/documents/vat-i_const_18700424_dei-filius_la.html}{official Latin text}.
\item Third Council of Constantinople (680--681), exposition of faith, DS 556; CCC 475: Christ's natural human will is neither suppressed nor opposed to his divine will but freely adheres to it in obedience. Council of Trent, sess.~VI, Decree on Justification (1547), chs.~1 and 5; canons 4--6, 17, DS 1521, 1525, 1554--1556, 1567: the fallen will is wounded rather than extinguished; prevenient grace and real cooperation. Canon 6 expressly condemns making evil works God's proper and per-se production rather than permission, as though Judas's betrayal were his work in the same manner as Paul's vocation. Trent sess.~XIII, ch.~4, DS 1642, governs the article's brief statement of transubstantiation.
\item Second Vatican Council, pastoral constitution \emph{Gaudium et spes} (1965), no.~17: authentic freedom as knowing choice of the good rather than license, healed by grace; \sourceurl{https://www.vatican.va/archive/hist_councils/ii_vatican_council/documents/vat-ii_const_19651207_gaudium-et-spes_en.html}{official English text}.
\item \emph{Catechism of the Catholic Church}, nos.~268--278 (universal, loving, wise, and mysterious omnipotence), \sourceurl{https://www.vatican.va/content/catechism/en/part_one/section_two/chapter_one/article_1/paragraph_3_the_almighty.html}{text}; 295, 301--324 (creation, concrete and immediate Providence, secondary causes, and the good God causes from permitted evil), \sourceurl{https://www.vatican.va/content/catechism/en/part_one/section_two/chapter_one/article_1/paragraph_4_the_creator.html}{text}; 599--600 and 606--614 (the Passion, divine plan, Christ's self-offering, and free agents), \sourceurl{https://www.vatican.va/content/catechism/en/part_one/section_two/chapter_two/article_4/paragraph_2_jesus_died_crucified.html}{text}; 1376 (transubstantiation); 1730--1748 (freedom and responsibility), \sourceurl{https://www.vatican.va/content/catechism/en/part_three/section_one/chapter_one/article_3.html}{text}; 1750--1761 (sources of morality); 1806 (prudence); 1857--1861 (grave matter, knowledge, consent, and judgment of persons); 1989--2002 (grace and freedom); 2725, 2734--2745 (effort, filial trust, and persevering prayer).
\item John Paul II, encyclical \emph{Veritatis splendor} (1993), nos.~34--41, 78--83: freedom participates in divine wisdom; intention does not erase the moral object; \sourceurl{https://www.vatican.va/content/john-paul-ii/en/encyclicals/documents/hf_jp-ii_enc_06081993_veritatis-splendor.html}{official English text}.
\item Innocent XI, apostolic constitution \emph{Caelestis Pastor} (1687), especially condemned propositions 2, 7--8, 12--14, 17, 31, 35, 37, and 40 (DS 2202, 2207--2208, 2212--2214, 2217, 2231, 2235, 2237, 2240) concerning inert passivity, concern for salvation and virtue, petition, temptation, and deliberate acts. Used to bound surrender against Quietism; \sourceurl{https://www.catholicculture.org/culture/library/view.cfm?recnum=5432}{working English text}.
\item Benedict XVI, First Vespers of Saints Peter and Paul (28 June 2008), on Augustine's maxim, a will united to Christ, truth, freedom, and responsibility; \sourceurl{https://www.vatican.va/content/benedict-xvi/en/homilies/2008/documents/hf_ben-xvi_hom_20080628_vespri.html}{official English text}.
\end{itemize}}

\subsection*{Patristic witnesses}

{\footnotesize
\begin{itemize}[itemsep=0.1em,topsep=0.15em]
\item St.~Irenaeus of Lyons, \emph{Against Heresies} \sourceurl{https://www.newadvent.org/fathers/0103404.htm}{IV.4.3} and \sourceurl{https://www.newadvent.org/fathers/0103437.htm}{IV.37.1--7}, on rational creatures as free, voluntary obedience, exhortation, praise, and blame; working English texts.
\item St.~Gregory of Nyssa, \emph{Great Catechism} 5--8, on the divine image, self-rule, and evil as recession or privation; \sourceurl{https://www.newadvent.org/fathers/29082.htm}{working English text}.
\item St.~John Chrysostom, \emph{Homilies on Matthew} 22.1--4, on Matt 6:25--34, work, care, and freedom from anxious servitude; \sourceurl{https://www.newadvent.org/fathers/200122.htm}{working English text}.
\item St.~John Damascene, \emph{Exposition of the Orthodox Faith} II.24--30, on voluntary and involuntary acts, choice and execution, Providence, permission, foreknowledge, and non-compulsion; \sourceurl{https://www.newadvent.org/fathers/33042.htm}{working English text}.
\item St.~Gregory the Great, \emph{Moralia in Iob} II.18.31, on Job's confession of divine governance through the narrated secondary causes. The image is used as spiritual exegesis, not to erase the raiders' or accuser's agency; \sourceurl{https://www.lectionarycentral.com/GregoryMoralia/Book02.html}{working English text}.
\item St.~Augustine, \emph{In epistulam Ioannis ad Parthos tractatus} VII.7--11, especially VII.8, \latin{Dilige, et quod vis fac}, PL 35, col.~2033; \sourceurl{https://www.newadvent.org/fathers/170207.htm}{working English text}. The passage's 1 John 4 context and moral contrasts control its interpretation.
\item Augustine, \emph{City of God} V.9--10 and XII.7--9 (foreknown order of causes, human wills, and deficient evil), \sourceurl{https://www.newadvent.org/fathers/120105.htm}{Book V} and \sourceurl{https://www.newadvent.org/fathers/120112.htm}{Book XII}; \emph{On Grace and Free Will} 2--5, 31--33 (commands, grace, willing, and cooperation), \sourceurl{https://www.newadvent.org/fathers/1510.htm}{text}; \emph{Enchiridion} 11, 95--102, 104--106 (omnipotent willed permission, the undefeated divine purpose, and freedom perfected), \sourceurl{https://www.newadvent.org/fathers/1302.htm}{text}; \emph{On Two Souls} 10--11 (force, prior fault, intention, and guilt), \sourceurl{https://www.newadvent.org/fathers/1403.htm}{text}.
\end{itemize}}

\subsection*{Thomistic metaphysical and moral framework}

{\footnotesize
\begin{itemize}[itemsep=0.1em,topsep=0.15em]
\item St.~Thomas Aquinas, \emph{Summa theologiae} I, q.~2, a.~3 (first cause in an essentially ordered series); q.~14, aa.~8, 11 (divine knowledge as causal and extending to singulars); q.~19, aa.~1--12, especially aa.~4, 6, 8, 9 ad 3, and 11--12 (the one simple divine will, its infallible efficacy, contingent mode, willed permission, and signs of will); q.~22, aa.~1--4 (Providence and its execution through secondary causes); q.~23, aa.~5--6 (first and secondary causality, certainty and contingent mode); q.~25, a.~3 (omnipotence and the inability to fail or sin); q.~44, a.~1 (participated being and cause); q.~46, a.~2 ad 7 (\latin{per se} and \latin{per accidens} series); qq.~48--49 (privation, moral fault, natural corruption, penalty, and the limp analogy); q.~83, aa.~1, 3--4 (free judgment, reason, and will); q.~103, aa.~5--8 (the universal government no creature defeats); q.~105, aa.~4--5 (God's interior motion of the will and operation in every agent).
\item Aquinas, \emph{Summa theologiae} I--II, q.~6, aa.~1--8 (voluntary and involuntary acts); qq.~9--10, especially q.~10, a.~4 (self-motion and divine motion of will without coercion); qq.~18--21 (object, intention, circumstances, and imputability); qq.~75--79, especially q.~79, aa.~1--2 (sin, ignorance, passion, and God causing the act's actuality but not its defect); qq.~109--114 (grace, free consent, cooperation, and merit). \sourceurl{https://www.corpusthomisticum.org/iopera.html}{Corpus Thomisticum Latin works}; \sourceurl{https://www.newadvent.org/summa/}{working English index}.
\item Aquinas, \emph{Summa theologiae} II--II, q.~23, a.~8; q.~25, a.~1; q.~26, aa.~1--2, on charity's form, loving the neighbor in God, and the order of love.
\item Aquinas, \emph{Summa contra Gentiles} III.67--74, 88--95, especially 70 (one whole effect from first and secondary agents in different orders), 73--74 (freedom and contingency), and 95 (prayer within Providence); \emph{De potentia} q.~3, a.~7; \emph{Super librum De causis} lect.~1.
\item The essentially / accidentally ordered comparison is presented as a Thomistic synthesis. ``Linear'' and ``hierarchical'' are this article's explanatory labels and not quotations from Aquinas.
\end{itemize}}

\subsection*{Spiritual witness and historical comparison}

{\footnotesize
\begin{itemize}[itemsep=0.1em,topsep=0.15em]
\item Jean Baptiste Saint-Jure, S.J., and St.~Claude de la Colombière, S.J., \emph{Trustful Surrender to Divine Providence: The Secret of Peace and Happiness}, trans.~Paul Garvin (Charlotte: TAN Books, modern reprint; originally published in English in 1961). The compilation identifies Part I as extracts from Saint-Jure's \emph{The Knowledge and Love of Our Lord Jesus Christ} and Part II as material drawn from La Colombière. Used especially at Saint-Jure I.1--2 and III, and La Colombière I.1 and I.3; \sourceurl{https://tanbooks.com/products/books/trustful-surrender-to-divine-providence-the-secret-of-peace-and-happiness/}{publisher record}.
\item The volume bears a 1961 imprimatur. That judgment of freedom from doctrinal or moral error is not a magisterial canonization of every rhetorical phrase or a critical authentication of the compilation's genealogy. The article's technical qualifications are governed by the primary doctrinal sources above.
\item St.~Ignatius of Loyola, \emph{The Spiritual Exercises}, no.~234, First Point of the ``Contemplation to Gain Love'': the boxed \emph{Suscipe} reproduces without adaptation the English translated from the Spanish autograph by Father Elder Mullan, S.J., \emph{The Spiritual Exercises of St.~Ignatius of Loyola} (New York: P.~J. Kenedy \& Sons, 1914), printed p.~120, \sourceurl{https://archive.org/details/spiritualexercis0000fath}{facsimile PDF p.~142}. The edition records the Jesuit printing faculty, Remigius Lafort's \emph{Nihil Obstat}, and John Cardinal Farley's \emph{Imprimatur} of 25 April 1914; it also reproduces an expressly abridged English text of Paul III's 31 July 1548 brief approving and praising the \emph{Exercises}. The papal brief concerns the \emph{Exercises}; the 1914 permissions concern Mullan's edition. Neither makes this a liturgical formula or a modern official translation.
\item Mullan's underlying 1914 United States text is public domain by publication date and remains received third-party material outside the project's CC BY 4.0 grant. A \sourceurl{https://ccel.org/ccel/i/ignatius/exercises/cache/exercises.pdf}{CCEL re-typeset PDF} was used only as a searchable cross-check; CCEL asserts separate rights in that file. This article incorporates neither verification PDF's wrapper, images, layout, pagination, nor typesetting.
\item Aleister Crowley, \emph{The Book of the Law: Liber AL vel Legis} (York Beach, ME: Red Wheel/Weiser, 2004), I.40, I.57, III.60; \sourceurl{https://lib.oto-usa.org/libri/liber0220.html}{O.T.O. online text and manuscript links}. Crowley's ``Duty'' was also checked in the \sourceurl{https://lib.oto-usa.org/crowley/essays/duty.html}{O.T.O. manuscript transcription}. These are authoritative evidence for Crowley's formulations, not neutral interpretations of them.
\item Crowley, ``The Antecedents of Thelema'' (1926), in \emph{The Revival of Magick and Other Essays}, ed.~Hymenaeus Beta and Richard Kaczynski (Tempe, AZ: New Falcon / O.T.O., 1998), checked against an \sourceurl{https://hermetic.com/eidolons/the-antecedents-of-thelema}{unfinished online transcript}, for his explicit retrospective acknowledgment of Augustine and Rabelais, his denial that Augustine's contextual meaning is that of \emph{Liber AL}, and his assignment of greater importance to Rabelais; Graham John Wheeler, ``Do What Thou Wilt: The History of a Precept,'' \emph{Religio} 27, no.~1 (2019): 17--41, especially 17--22, \sourceurl{https://hdl.handle.net/11222.digilib/141542}{stable record and PDF}, for the multiple genealogy and historical qualification. These sources warrant comparison but do not establish direct borrowing from Augustine in 1904; the article's limited theological contrast is editorial synthesis, not a claim of textual descent or a taxonomy of movements.
\end{itemize}}

\endgroup
